% 5.6 — Limites assumées
\section{Limites assumées}
\label{sec:limites}

Un garde-fou qui cache ses limites crée une fausse confiance. Celles de \demo{} sont
connues, et la plupart sont dites dans la raison même des décisions.

\begin{enumerate}
  \item \textbf{Des erreurs, pas un adversaire.} Le modèle de menace vise un agent qui se trompe.
    Un adversaire déterminé dispose de moyens que l'outil ne prétend pas couvrir (charge
    chiffrée, course entre vérification et usage d'un appel système).
  \item \textbf{Le filet Linux est un prototype.} ptrace sur x86-64 seulement, un arrêt par appel
    système ; la cible de production désignée par le code est seccomp avec notification.
  \item \textbf{Les filets sauvegardent le projet, et seulement lui.} Hors du projet, une
    destruction par du code non lu n'est pas sauvegardée ; le filet Linux l'enregistre (un reçu
    par fichier), le montage Windows ne la voit pas.
  \item \textbf{Un exécuteur de code sans filet est autorisé} avec une note qui dit qu'aucune
    capture n'a été faite (arbitrage M1) --- sauf le code en ligne qui lance d'autres programmes
    et les installations de paquets, qui escaladent (Rb2).
  \item \textbf{Un mot sans règle n'a pas d'effet} (Rb1). Un outil absent de la table et qui n'est
    pas un exécuteur de code est traité comme sans effet ; la table s'étend famille par famille
    pour réduire ce cas.
  \item \textbf{Ce que C3 ne lit pas.} Une construction PowerShell hors de la grammaire de C3
    (chaîne construite à l'exécution, classe, fonction) est non comprise : filet WinFsp, sinon
    point de contrôle, sinon escalade ; un script \cmd{.ps1} est un exécuteur de code.
  \item \textbf{Codex ne sait que refuser} : pas de question, pas de filet par réécriture (V15).
  \item \textbf{Le distant ne se récupère pas.} Pour une base, un dépôt ou un service distant,
    \demo{} reconnaît, demande ou refuse, et indique où le fournisseur garde son historique ;
    il ne restaure rien lui-même.
  \item \textbf{Les points de contrôle sont optionnels} et bornés ; désactivés, une ligne non lue
    sans filet est escaladée.
\end{enumerate}

\section*{Conclusion du chapitre}
La sécurité de \demo{} n'a pas été supposée : elle a été attaquée --- par la lecture, par des
sondes, par la fouille des registres --- et chaque faille trouvée a laissé un test. La
validation repose sur deux filets complémentaires : les tests vérifient ce qu'on a pensé à
écrire, le portail montre tout ce qui bouge. Les limites qui restent sont connues, datées et
planifiées.
